% !TEX root = task-1-17-2.tex
% Editable task bank. Append new BankTask blocks at the end of the document.
\documentclass[12pt,a4paper]{article}
\usepackage{physics-textbook}
\usepackage{graphicx,needspace}
\usetikzlibrary{decorations.pathmorphing}
\newif\ifShowSolutions
\ifdefined\TutorEdition
  \ShowSolutionstrue
\else
  \ShowSolutionsfalse % Change to \ShowSolutionstrue to display solutions.
\fi
\newcounter{banktask}
\newsavebox{\BankTaskBox}
\newcommand{\BankTaskContent}[3]{%
  \noindent\textbf{\thebanktask.}\quad #2\par
  \ifShowSolutions
    \Needspace{50mm}% Keep the solution heading with its opening text or diagram.
    #3\par
    \ifdefined\TutorEdition
      \ifcsname TutorSource#1\endcsname
        \nopagebreak\smallskip{\small\csname TutorSource#1\endcsname\par}%
      \fi
    \fi
  \fi
}
\newcommand{\BankTask}[3]{%
  \refstepcounter{banktask}\par\medskip
  \sbox{\BankTaskBox}{\begin{minipage}{\linewidth}%
    \BankTaskContent{#1}{#2}{#3}\end{minipage}}%
  \ifdim\dimexpr\ht\BankTaskBox+\dp\BankTaskBox\relax>\dimexpr\textheight-20mm\relax
    \Needspace{35mm}\BankTaskContent{#1}{#2}{#3}%
  \else
    \Needspace{\dimexpr\ht\BankTaskBox+\dp\BankTaskBox\relax}%
    \noindent\usebox{\BankTaskBox}\par
  \fi
  \medskip
}
\begin{document}
\BookSubsection{1.17.2}{Уравнение гармонических колебаний}


% Source sdamgia-635: https://phys-ege.sdamgia.ru/problem?id=635
\BankTask{sdamgia-635}{
Колебательное движение тела задано уравнением:

$x=a \sin ( bt + \frac{\pi}{2} ) ,$

где $a = 5\,\text{см},$ $b=3\,\text{с} ^{ - 1 } .$ Чему равна амплитуда колебаний?
}{
\smallskip\noindent\textbf{Решение}\par
Общий вид закона изменения со временем координаты тела, совершающего колебания, имеет вид:

$x=x_0 \sin ( \omega t + \varphi_0 ) ,$

где $x_0$  — амплитуда колебаний.

Сравнивая с $x=a \sin ( bt + \frac{\pi}{2} ) ,$ заключаем, что амплитуда колебаний равна $x_0=a=5\,\text{см}.$

Ответ: $5$.
}

% Source sdamgia-6113: https://phys-ege.sdamgia.ru/problem?id=6113
\BankTask{sdamgia-6113}{
Два маятника $1$ и $2$ совершают гармонические колебания по законам $x_1(t)=9\sin2t$ и $x_2(t)=3\cos2t$. Фазы колебаний этих маятников

\noindent 1) одинаковые\par
\noindent 2) отличаются в $3$ раза\par
\noindent 3) отличаются в $4{,}5$ раза\par
\noindent 4) отличаются на $\pi /2$\par
}{
\smallskip\noindent\textbf{Решение}\par
Воспользуемся формулами приведения:

$x_2=3 \cos 2t=3 \sin ( 2t + \frac{\pi}{2} ) .$

Следовательно фазы колебаний этих маятников отличаются на $\frac{\pi}{2} .$

Правильный ответ указан под номером $4$.
}

% Source sdamgia-8429: https://phys-ege.sdamgia.ru/problem?id=8429
\BankTask{sdamgia-8429}{
Гиря массой $2\,\text{кг}$ подвешена на стальной пружине и совершает свободные колебания вдоль вертикально направленной оси $Ox$, координата $x$ центра масс гири, выраженная в метрах, изменяется со временем по закону $x = 0{,}4 \cdot \sin 5t.$ Чему равна кинетическая энергия гири в начальный момент времени?
}{
\smallskip\noindent\textbf{Решение}\par
Поскольку координата центра масс гири меняется по закону:

$x=0{,}4 \cdot \sin 5t = x_0 \sin \omega_0t,$

скорость центра масс гири меняется согласно:

$v =x_0\omega_0 \cos \omega_0t=0{,}4 \cdot 5 \cdot \cos 5t=2 \cdot \cos 5t.$

В начальный момент скорость гири равна $2\,\text{м/с}$, а ее кинетическая энергия:

$E= \frac{m v ^2}{2} = \frac{2 \cdot 2 ^2}{2} =4\,\text{Дж}.$

Ответ: $4$.
}

% Source sdamgia-720: https://phys-ege.sdamgia.ru/problem?id=720
\BankTask{sdamgia-720}{
Гиря массой $2\,\text{кг}$ подвешена на стальной пружине и совершает свободные колебания вдоль вертикально направленной оси $Ox$, координата $x$ центра масс гири изменяется со временем по закону $x = 0{,}4 \cdot \sin 5t.$ Кинетическая энергия гири изменяется по закону

\noindent 1) $4 \cdot \cos ^2 5t$ \par
\noindent 2) $8 \cdot \sin ^2 5t$ \par
\noindent 3) $4 \cdot \sin ^2 5t$ \par
\noindent 4) $8 \cdot \cos ^2 5t$\par
}{
\smallskip\noindent\textbf{Решение}\par
Поскольку координата центра масс гири меняется по закону 

$x=0{,}4 \cdot \sin 5t = x_0 \sin \omega_0t,$

скорость центра масс гири меняется согласно

$v =x_0\omega_0 \cos \omega_0t=0{,}4 \cdot 5 \cdot \cos 5t=2 \cdot \cos 5t.$

Следовательно, кинетическая энергия изменяется по закону

$E= \frac{m v ^2}{2} = \frac{2 \cdot ( 2 \cdot \cos 5t ) ^2}{2} =4 \cdot \cos ^2 5t.$

Ответ: $1$.
}

% Source sdamgia-10063: https://phys-ege.sdamgia.ru/problem?id=10063
\BankTask{sdamgia-10063}{
Смещение груза горизонтального пружинного маятника от положения равновесия меняется с течением времени по закону $x = A \cos ( \frac{2 \pi}{T} t ) ,$ где период $T$  =  $1\,\text{с}$. Через какое минимальное время начиная с момента $t$  =  $0$ потенциальная энергия маятника вернется к своему исходному значению?
}{
\smallskip\noindent\textbf{Решение}\par
Потенциальная энергия пружинного маятника зависит от квадрата смещения груза:

$W= \frac{kx ^2}{2} .$

В нулевой момент времени смещение груза $x_1=A.$ Через $t = \frac{T}{2} =0{,}5\,\text{с}$ смещение груза будет равно:

$x=A \cos ( \frac{2 \pi}{T} \cdot \frac{T}{2} ) = - A.$

И потенциальная энергия вернется к своему исходному значению.

Ответ: $0{,}5$.
}

% Source sdamgia-25358: https://phys-ege.sdamgia.ru/problem?id=25358
\BankTask{sdamgia-25358}{
Изменение координаты при движении маятника имеет вид $x=A \sin ( \frac{2 \pi}{T} t ) .$ Период колебаний маятника равен $T$  =  $2\,\text{с}$. Через какое время кинетическая энергия маятника впервые примет минимальное значение?
}{
\smallskip\noindent\textbf{Решение}\par
Так как уравнение зависимости $x(t)$ выражено через синус аргумента, то можно сделать вывод, что в начальный момент времени тело находилось в положении равновесия. Кинетическая энергия минимальна в крайних точках, в которых потенциальная энергия максимальна, то есть смещение тела равно амплитуде $x=A$. Значит, $\sin ( \frac{2 \pi}{T} t ) =1.$ Тогда: 

$\frac{2 \pi}{T} t= \frac{\pi}{2} \Longleftrightarrow t= \frac{T}{4} = \frac{2}{4} =0{,}5$ с.

Ответ: $0{,}5$.
}

% Source sdamgia-29061: https://phys-ege.sdamgia.ru/problem?id=29061
\BankTask{sdamgia-29061}{
Проекция скорости колеблющегося тела изменяется по закону $v _{x} = 0{,}5 \cos ( 4t ) .$ Чему равен модуль ускорения тела в момент, когда проекция скорости равна $30\frac{\,\text{см}}{\,\text{с}}$?
}{
\smallskip\noindent\textbf{Решение}\par
$1$.  Проекция ускорения является производной от скорости по времени. Тогда уравнение изменения ускорения будет иметь вид:

$a_x= v '_t= - 0{,}5 \cdot 4 \sin ( 4t ) = - 2 \sin ( 4t ) .$ 

$2$.  Проекция скорости тела будет равна $30\,\text{см/с}$  =  $0{,}3\,\text{м/с}$ при условии $0{,}3=0{,}5 \cos ( 4t ) ,$ откуда $\cos ( 4t ) =0{,}6.$ Тогда: 

$\sin(4t)=\sqrt{1-\cos^2(4t)}=0{,}8.$ 

$3$.  Найдем проекцию ускорения в этот момент времени: 

$a_x= - 2 \cdot 0{,}8= - 1{,}6\; \frac{\,\text{м}}{\,\text{с}^{2}} .$ 

Модуль ускорения $a = 1{,}6 \; \frac{\,\text{м}}{\,\text{с}^{2}} .$

Ответ: $a$ = $1{,}6\,\text{м}/\text{с}^{2}$.
}

% Source sdamgia-52049: https://phys-ege.sdamgia.ru/problem?id=52049
\BankTask{sdamgia-52049}{
Небольшой груз массой $m$  =  $200\,\text{г}$ совершает незатухающие гармонические колебания на невесомой нерастяжимой нити длиной $l$  =  $2\,\text{м}$. Закон движения этого груза имеет вид: $x ( t ) = 0{,}2 \cos ( \omega t ) ,$ где величина смещения $x$ от положения равновесия измеряется в метрах. Определите максимальную кинетическую энергию груза в процессе колебаний.
}{
\smallskip\noindent\textbf{Решение}\par
Максимальная кинетическая энергия груза равна: 

$E_k = \frac{m v _{\max} ^2}{2} , \qquad ( 1 )$

где максимальная скорость груза в процессе колебаний 

$v _{\max} = \omega A, \qquad ( 2 )$

а $A$  — амплитуда колебаний маятника, величину которой можно найти из закона движения: $A = 0{,}2 \,\text{м}.$
Частота $\omega$ и период $T$ колебаний маятника связаны соотношением: 

$\omega = \frac{2 \pi}{T} . \qquad ( 3 )$

Период колебаний математического маятника можно найти по формуле: 

$T = 2 \pi \sqrt{\frac{l}{g}} . \qquad ( 4 )$

Тогда из уравнений ($1$)–($4$) получим: 

$E_k = \frac{m g A ^2}{2 l} = \frac{0{,}2 \cdot 10 \cdot 0{,}2 ^2}{2 \cdot 2} = 0{,}02 \text{Дж}.$

Ответ: $0{,}02\,\text{Дж}$.
}


% Source eeB001: https://ege.fipi.ru/bank/questions.php?search=1&proj=BA1F39653304A5B041B656915DC36B38&qid=eeB001
\BankTask{eeB001}{
В таблице представлены данные о положении грузика, прикреплённого к пружине и совершающего гармонические колебания вдоль горизонтальной оси $Ox$, в различные моменты времени.
\begin{center}{\small\setlength{\tabcolsep}{3pt}\begin{tabular}{|c|c|c|c|c|c|c|c|c|c|c|c|c|c|c|c|c|}\hline $t,\,\text{с}$ & $0{,}0$ & $0{,}1$ & $0{,}2$ & $0{,}3$ & $0{,}4$ & $0{,}5$ & $0{,}6$ & $0{,}7$ & $0{,}8$ & $0{,}9$ & $1{,}0$ & $1{,}1$ & $1{,}2$ & $1{,}3$ & $1{,}4$ & $1{,}5$ \\\hline $x,\,\text{мм}$ & $0$ & $2{,}5$ & $4{,}5$ & $6{,}0$ & $7{,}0$ & $7{,}5$ & $7{,}0$ & $6{,}0$ & $4{,}5$ & $2{,}5$ & $0$ & $-2{,}5$ & $-4{,}5$ & $-6{,}0$ & $-7{,}0$ & $-7{,}5$ \\\hline\end{tabular}}\end{center}
\begin{enumerate}
\item Период колебаний грузика равен $1\,\text{с}$.
\item Амплитуда колебаний грузика равна $15\,\text{мм}$.
\item Потенциальная энергия пружины маятника в момент времени $1{,}5\,\text{с}$ минимальна.
\item Кинетическая энергия грузика в момент времени $1{,}0\,\text{с}$ максимальна.
\item Модуль ускорения грузика в момент времени $1{,}5\,\text{с}$ максимален.
\end{enumerate}
}{
\smallskip\noindent\textbf{Решение от GPT}\par
Из таблицы $A=7{,}5\,\text{мм}$, $T=2\,\text{с}$. При $t=1{,}0\,\text{с}$ $x=0$, поэтому кинетическая энергия максимальна; при $t=1{,}5\,\text{с}$ $|x|=A$, поэтому $|a|=\omega^2A$ максимально. Ответ: $45$.
}

% Source 9F061F: https://ege.fipi.ru/bank/questions.php?search=1&proj=BA1F39653304A5B041B656915DC36B38&qid=9F061F
\BankTask{9F061F}{
Тело совершает гармонические колебания вдоль оси $Ox$. В таблице приведены координаты этого тела через равные промежутки времени.
\begin{center}{\small\setlength{\tabcolsep}{3pt}\begin{tabular}{|c|c|c|c|c|c|c|c|}\hline $t,\,\text{с}$ & $0$ & $0{,}2$ & $0{,}4$ & $0{,}6$ & $0{,}8$ & $1{,}0$ & $1{,}2$ \\\hline $x,\,\text{см}$ & $0$ & $2$ & $4$ & $2$ & $0$ & $2$ & $4$ \\\hline\end{tabular}}\end{center}
\begin{enumerate}
\item Амплитуда колебаний тела равна $8\,\text{см}$.
\item Частота колебаний тела равна $1{,}25\,\text{Гц}$.
\item Модуль ускорения тела максимален в момент времени $0{,}2\,\text{с}$.
\item В момент времени $0{,}4\,\text{с}$ тело проходит положение равновесия.
\item В момент времени $0{,}6\,\text{с}$ тело обладает максимальной кинетической энергией.
\end{enumerate}
}{
\smallskip\noindent\textbf{Решение от GPT}\par
Положение равновесия $x_{\text{равн}}=2\,\text{см}$, амплитуда $A=2\,\text{см}$, период $T=0{,}8\,\text{с}$; $\nu=1/T=1{,}25\,\text{Гц}$. В момент $0{,}6\,\text{с}$ тело в равновесии и его кинетическая энергия максимальна. Ответ: $25$.
}

% Source C1C6AE: https://ege.fipi.ru/bank/questions.php?search=1&proj=BA1F39653304A5B041B656915DC36B38&qid=C1C6AE
\BankTask{C1C6AE}{
На рисунке приведены зависимости координат $x$ двух грузов $1$ и $2$ от времени $t$ в процессе их малых колебаний на гладкой горизонтальной плоскости. Грузы имеют одинаковую массу $m$ и прикреплены к пружинам жёсткостью $k_{1}$ (груз $1$) и жёсткостью $k_{2}$ (груз $2$). Колебания груза $1$ обозначены сплошной линией, груза $2$ – пунктиром.
\begin{center}\begin{tikzpicture}[x=.96cm,y=.8cm,>=stealth]
\foreach \q in {0,1,2,3,4,5,6} \draw[thin,black!65] (\q,-2)--(\q,2);
\draw[thin,black!65] (0,2)--(6,2) (0,-2)--(6,-2);
\foreach \h in {-1.5,-1,-.5,.5,1,1.5} \draw[thin,black!65] (0,\h)--(1.1,\h) (1.4,\h)--(2,\h) (2.8,\h)--(3.5,\h) (4,\h)--(5,\h) (5.3,\h)--(6,\h);
\draw[->] (0,0)--(6.6,0) node[right] {$t,\,\text{с}$};
\draw[->] (0,-2)--(0,2.7) node[above] {$x,\,\text{см}$};
\foreach \x in {1,2,3,4,5,6} \node[below] at (\x,0) {$\x$};
\foreach \y in {-2,-1,0,1,2} \node[left] at (0,\y) {$\y$};
\draw[very thick,domain=0:6,samples=100] plot (\x,{2*sin(90*\x)});
\draw[thick,dashed,domain=0:6,samples=100] plot (\x,{cos(90*(\x-.25))});
\node at (5.8,1.85) {$1$};\node at (3.2,1.0) {$2$};
\end{tikzpicture}\end{center}
\begin{enumerate}
\item Периоды колебаний грузов одинаковы.
\item Максимальные потенциальные энергии деформированных пружин одинаковы.
\item Жёсткость пружины $k_{1}$ меньше жёсткости пружины $k_{2}$.
\item Амплитуды колебаний грузов одинаковы.
\item Максимальная кинетическая энергия груза $1$ в четыре раза больше, чем максимальная кинетическая энергия груза $2$.
\end{enumerate}
}{
\smallskip\noindent\textbf{Решение от GPT}\par
По графику периоды одинаковы, значит при одинаковых массах $k_1=k_2$. Амплитуды $A_1=2\,\text{см}$ и $A_2=1\,\text{см}$ различаются вдвое. Поэтому $E_{\text{к max}}=kA^2/2$ у первого груза вчетверо больше. Ответ: $15$.
}

% Source BE8E13: https://ege.fipi.ru/bank/questions.php?search=1&proj=BA1F39653304A5B041B656915DC36B38&qid=BE8E13
\BankTask{BE8E13}{
Один конец лёгкой пружины жёсткостью $k$ прикреплён к бруску, а другой закреплён неподвижно. Брусок скользит по горизонтальной направляющей так, что его координата изменяется со временем по закону $x(t)=A\sin\omega t$.
Установите соответствие между физическими величинами и формулами, выражающими их изменения во времени.\TaskMatchingColumns[.53\linewidth]{Физические величины}{
А) проекция $F_x(t)$ равнодействующей силы на ось $x$\par
Б) потенциальная энергия пружины $E_{\text{П}}(t)$
}{Формулы}{
1) $kA^2\sin\omega t$\par
2) $-kA\sin\omega t$\par
3) $-kA^2\sin^2\omega t$\par
4) $\frac{kA^2}{2}\sin^2\omega t$\par
}\begin{center}\begin{tabular}{|c|c|}\hline А & Б \\\hline\rule{0pt}{1.2em}\hspace{1cm} & \hspace{1cm} \\\hline\end{tabular}\end{center}
}{
\smallskip\noindent\textbf{Решение от GPT}\par
По закону Гука $F_x=-kx=-kA\sin\omega t$ (формула $2$). Энергия $E_{\text{П}}=kx^2/2=(kA^2/2)\sin^2\omega t$ (формула $4$). Ответ: $24$.
}

% Source 731: https://phys-ege.sdamgia.ru/problem?id=731
\BankTask{731}{
Подвешенный на нити грузик совершает гармонические колебания. В таблице представлены координаты грузика через одинаковые промежутки времени. Какова примерно максимальная скорость грузика? 

\begin{center}\begin{tabular}{|c|c|c|c|c|c|c|c|c|}\hline
t, c&$0$&$0{,}5$&$1$&$1{,}5$&$2$&$2{,}5$&$3$&$3{,}5$\\\hline
x, см&$6$&$3$&$0$&$3$&$6$&$3$&$0$&$3$\\\hline
\end{tabular}\end{center}
}{
\smallskip\noindent\textbf{Решение}\par
Максимальная скорость груза маятника связана с амплитудой колебаний и циклической частотой соотношением: 

\[v _{\,\text{м}}=x_{\,\text{м}} \omega_0.\]

Из таблицы видно, что при колебаниях смещение меняется от $0$ до $6\,\text{см}$. Следовательно, в положении равновесия смещение равно $3\,\text{см}$, а амплитуда колебаний равна $x_{\,\text{м}}=3\,\text{см},$ период $T=2\,\text{с},$ а значит, циклическая частота:

\[\omega_0= \frac{2 \pi}{T} = \frac{2 \pi}{2\,\text{с}} = \pi \,\text{рад/с}.\]

Таким образом, максимальная скорость груза равна приблизительно:

\[v _{\,\text{м}}=0{,}03\,\text{м} \cdot 3{,}14\,\text{с}^{-1}\approx 0{,}09\,\text{м/с}.\]

Ответ: $0{,}09\,\text{м/с}$.
}

\end{document}
